\bookchapter{Capítulo 26: Fronteras tecnológicas actuales del móvil}

\chapterintro{Tras hablar de la “despedida” de los teléfonos antiguos, volvamos al presente. 2024-2026 es un periodo de explosión tecnológica: la IA se ha colocado en el bolsillo desde la nube, la comunicación por satélite ha bajado de los gama alta a los gama media, las pantallas plegables han pasado de “novedad” a “uso diario”, y la tecnología de baterías ha visto un avance a nivel de materiales no visto en muchos años. En este capítulo repasamos estas tecnologías emergentes para ayudarte a decidir “si ahora es el momento de cambiar” y “a qué cambiar”.}

\booksection{I. ¿AI phone: es un truco de marketing o un cambio real?}

Desde 2024, casi todos los fabricantes Android han lanzado teléfonos insignia bajo la bandera de “AI phone”, y Apple también lanzó Apple Intelligence en iOS 18. Pero, ¿qué tipo de teléfono cuenta realmente como “AI phone”? La industria ha formado aproximadamente tres criterios:

\begin{itemize}
  \item \textbf{Grandes modelos en el dispositivo}: El teléfono puede ejecutar localmente modelos de lenguaje de unos 7B de parámetros (como versiones ligeras de Qwen, Llama, DeepSeek), sin depender de la nube para completar resúmenes de texto, traducción, generación de textos, etc. Método de verificación: si las funciones de IA siguen funcionando en modo avión, es un modelo en el dispositivo.
  \item \textbf{Capacidad de NPU (TOPS)}: La capacidad de la NPU (unidad de procesamiento neuronal) se mide en TOPS (billones de operaciones por segundo). Los chips insignia de 2024 (como Snapdragon 8 Gen3, Dimensity 9300) tienen una capacidad de NPU entre 40-70 TOPS; los nuevos chips de 2025-2026 superangeneralmentemente los 100 TOPS, este es el umbral de hardware para ejecutar grandes modelos en el dispositivo.
  \item \textbf{Escenarios de AI Agent}: No solo "responde preguntas", sino que puede ejecutar cadenas de tareas entre apps; por ejemplo, "organiza las fotos de la reunión de ayer en un álbum y compártelo con Zhang San", el sistema puede invocar automáticamente las apps de galería, archivos, WeChat, etc. Esta es la marca distintiva de un AI phone frente a un "teléfono con asistente de IA".
\end{itemize}

\booksection{II. Comparación de estrategias de IA: ¿quién lo hace en serio y quién solo pone una etiqueta?}

Las rutas de IA de cada fabricante son claramente diferentes; al comprar puedes prestar atención a:

\begin{itemize}
  \item \textbf{Apple Apple Intelligence (iOS 18+)}: Se centra en la integración de apps del sistema: resúmenes de correo, prioridad de notificaciones, corrección de escritura, edición de fotos con IA. Principalmente modelos en el dispositivo; las tareas complejas van a la nube con Private Cloud Compute. La ventaja es la protección de privacidad y la fluidez del sistema; la desventaja es que las funciones en chino llegan más lento, algunas funciones no se soportan hasta 2025.
  \item \textbf{Google Gemini Nano}: Integrado en los Pixel 8/9, soporta "Círculo para buscar", Magic Compose, resumen de grabaciones, etc. Gemini Nano es un pequeño modelo en el dispositivo en el sentido más auténtico; el ecosistema de código abierto es amigable y es la implementación de referencia para los AI phones de Android.
  \item \textbf{." (vida cotidiana)}: Integrado en HarmonyOS Xiaoyi, soporta generación de textos, edición de fotos con IA, eliminación con IA (quitar personas), resumen de llamadas. Tras HarmonyOS NEXT 5.0, la integración a nivel de sistema es mayor, las capacidades de IA se fusionan profundamente con el sistema.
  \item \textbf{Modelo grande MiLM de Xiaomi}: Integrado en HyperOS, soporta retratos con IA, resúmenes con IA, pintor con IA, superresolución. Xiaomi sigue una ruta de "colaboración nube-dispositivo": modelos ligeros en el dispositivo para tareas pesadas en la nube.
  \item \textbf{OPPO (modelo grande Andes) y vivo (modelo grande Blue Heart)}: Integrados en ColorOS 15 y OriginOS 5, se centran en eliminación con IA, resumen de documentos con IA, preguntas con reconocimiento de imágenes, búsqueda de fotos en lenguaje natural ("buscar fotos tomadas en la playa el año pasado").
\end{itemize}

\textbf{Recomendacion de compra}: Si necesitas frecuentemente IA para procesar documentos de trabajo y organizar información, prioriza modelos con gran capacidad de cómputo en el dispositivo e integración de IA a nivel de sistema; si solo usas IA para editar fotos ocasionalmente, los gama media ya son suficientes.

\booksection{III. Comunicación por satélite: de “truco” de gama alta a estándar de gama media}

La comunicación por satélite es una de las tecnologías que más ha bajado en los últimos dos años. En 2022, con el Huawei Mate 60 y el iPhone 14, era exclusiva de los gama alta; en 2025 ya ha bajado al rango de 2000-3000 yuanes.

\begin{itemize}
  \item \textbf{Mensajes cortos de BeiDou en gama media}: La función de mensajes cortos del satélite BeiDou (puede enviar SMS de hasta 1000 caracteres chinos + ubicación) se ha extendido a modelos gama media de Redmi, realme, iQOO, etc., sin cuota mensual, uso gratuito. Adecuado para senderismo al aire libre, viajes en coche, etc.
  \item \textbf{Llamadas por satélite Tiantong}: Los Huawei Mate 60/70 y Honor Magic6 soportan llamadas por satélite Tiantong-1, pudiendo hacer/recibir llamadas por satélite y enviar mensajes satelitales. Requiere activar un plan de satélite del operador (aprox. 10-200 yuanes/mes), adecuado para usuarios que entran frecuentemente en zonas sin cobertura.
  \item \textbf{Mensajes satelitales de iPhone (SOS)}: iPhone 14 y posteriores soportan SOS por satélite, pudiendo enviar mensajes de texto a servicios de emergencia y compartir la ubicación. Servicio gratuito durante 2 años, con renovación disponible en algunas regiones.
  \item \textbf{Consejos prácticos}: Los usuarios urbanos básicamente no lo necesitan, pero se recomienda encarecidamente a entusiastas de actividades al aire libre y profesionales de negocios que viajan a zonas remotas que elijan modelos con función satelital; puede salvar vidas en momentos críticos.
\end{itemize}

\booksection{IV. Madurez del ecosistema de pantallas plegables: ¿se puede usar como teléfono principal?}

Las pantallas plegables han pasado de “novedad tecnológica” a “aptas como teléfono principal”. Para evaluar la madurez de una pantalla plegable, hay tres aspectos: durabilidad del hardware, adaptación de apps, ecosistema de interacción.

\begin{itemize}
  \item \textbf{Durabilidad del hardware}: La vida útil de la bisagra alcanzageneralmentemente más de 500.000 pliegues (100 pliegues diarios durante 13 años), las bisagras de gota de agua reducen el pliegue de pantalla, la resistencia al agua IPX8 es estándar. La pantalla interior con vidrio UTG + película protectora hace que la pantalla ya no sea "delicada".
  \item \textbf{Adaptación de apps}: Las apps principales (WeChat, Taobao, Douyin, Bilibili, Office, WPS) se han adaptado a pantallas grandes, pero todavía hay un 30\% de apps de cola larga sin adaptar, mostrándose estiradas o con bordes negros. Se recomienda probar tus apps habituales antes de comprar.
  \item \textbf{Interacción en modo suspendido}: Al plegar entre 90°-120°, algunas apps soportan el modo suspendido donde "la mitad superior muestra el contenido y la inferior permite operar", adecuado para videollamadas, selfies y ver videos. La cámara con selfie en modo suspendido es una experiencia única de los plegables.
  \item \textbf{Pantalla dividida y multitarea}: En pantalla grande se pueden abrir dos apps en pantalla dividida e incluso tres apps en ventanas pequeñas simultáneamente. HyperOS, ColorOS y OriginOS ya tienen lógicas de pantalla dividida bastante maduras, pero iOS (estilo iPadOS) aún no permite pantalla dividida en el iPhone.
  \item \textbf{Precio}: Los plegables horizontales (tipo libro) han bajado del rango de 10.000+ a 5.000-7.000 yuanes; los plegables verticales (tipo concha) han bajado a 4.000-5.000 yuanes. En 2025, algunos fabricantes lanzaron plegables verticales en el rango de 3.000 yuanes, entrando verdaderamente en el mercado masivo.
\end{itemize}

\textbf{Recomendacion de compra}: Si eres un usuario intensivo de lectura, procesamiento de documentos y video, los plegables horizontales merecen consideración; si buscas portabilidad y atributos sociales, los plegables verticales son más adecuados. Se recomienda probar en tienda física al menos 30 minutos antes de decidir.

\booksection{V. 6G y Wi-Fi 7: velocidad y límites de la próxima generación de conectividad}

Aunque el 6G aún no se ha implementado realmente, las tecnologías relacionadas ya están en camino:

\begin{itemize}
  \item \textbf{Progreso del 6G}: 3GPP planea completar el estándar 6G en 2027, con comercialización hacia 2030. La velocidad pico teórica del 6G puede alcanzar 1 Tbps (50 veces la del 5G), con latencia tan baja como 0.1ms, soportando "integración de comunicación y percepción"; el teléfono podría incluso funcionar como un radar de baja precisión. Actualmente los fabricantes siguen en fase de prototipo, lejos de productos de consumo.
  \item \textbf{Wi-Fi 7(802.11be)}: Desde 2024 es estándar en teléfonos insignia; la velocidad pico teórica alcanza 46 Gbps (4 veces la del Wi-Fi 6), soporta ancho de banda de 320MHz, 4K-QAM y MLO (multi-enlace simultáneo). En la práctica, con un router Wi-Fi 7, los teléfonos insignia pueden alcanzar 2-5 Gbps, 2-3 veces más rápido que Wi-Fi 6E. Si compras un teléfono con Wi-Fi 7, recuerda combinarlo con un router Wi-Fi 7 para aprovechar todo el rendimiento.
  \item \textbf{Bluetooth 5.4 y LE Audio}: Los teléfonos nuevos de 2024-2026 soportangeneralmentemente Bluetooth 5.4 y LE Audio (audio de baja consumo), pudiendo conectar dos auriculares simultáneamente para escuchar la misma canción; Auracast permite que varias personas escuchen la misma fuente de audio en espacios públicos.
\end{itemize}

\booksection{VI. Batería de silicio-carbono: la autonomía por fin ve un avance a nivel de materiales}

La autonomía de los teléfonos ha sido criticada durante años por “una carga al día”; desde 2024, la popularización de la tecnología de batería de silicio-carbono ha mejorado la situación de forma sustancial.

\begin{itemize}
  \item \textbf{Principio técnico}: La capacidad teórica del ánodo de grafito tradicional es de unos 372 mAh/g; la del silicio alcanza los 4200 mAh/g (unas 10 veces), pero el silicio se expande un 300\% durante la carga/descarga, tendiendo a pulverizarse. El ánodo de silicio-carbono, mediante una estructura compuesta de nano-silicio y carbono, controla el problema de expansión del silicio dentro de un rango aceptable, aumentando la densidad energética de la batería en un 15\%-30\%.
  \item \textbf{realmente adaptadas a ti}: Los teléfonos insignia de 2024-2025 con batería de silicio-carbono (como Xiaomi 14 Ultra, Honor Magic 6 Ultimate, vivo X100 Pro, etc.) consiguen meter baterías de 5000-6000 mAh en cuerpos de 7-8mm, unas 800-1000 mAh más que la generación anterior; en uso intensivo pueden durar un día y medio.
  \item \textbf{Combinación con carga rápida de 100W+}: El ánodo de silicio-carbono es compatible con la carga rápida; 100W/120W que carga completamente en 17 minutos ya no es noticia. En 2025, algunos modelos soportan 150W-240W, cargando al 80\% en 10 minutos.
  \item \textbf{Recomendacion de compra}: Si eres un usuario intensivo, prioriza modelos posteriores a 2024 con la combinación "silicio-carbono + carga rápida de 100W+"; la ansiedad por la batería se reducirá drásticamente.
\end{itemize}

\booksection{VII. Resumen del capítulo: recuerda estas tendencias en una frase}

\begin{itemize}
  \item Tres criterios de AI phone: modelo grande en el dispositivo, capacidad NPU TOPS, escenarios de AI Agent entre apps.
  \item Comunicación por satélite: mensajes cortos de BeiDou en gama media, llamadas Tiantong siguen siendo de gama alta.
  \item Pantallas plegables: hardware maduro, adaptación de apps al 70\%, modo suspendido/pantalla dividida son experiencias únicas.
  \item El 6G está en camino, el Wi-Fi 7 ya está aquí, el LE Audio de Bluetooth empieza a popularizarse.
  \item Batería de silicio-carbono + carga rápida de 100W+: doble avance en autonomía y carga.
\end{itemize}

\begin{tcbpractice}{Pruébalo: haz una "lista de requisitos técnicos" para tu próximo teléfono}
Coge un papel y, en tres categorías “imprescindible/bonus/no necesario”, lista tus necesidades de IA, satélite, pantalla plegable, Wi-Fi 7, carga rápida y capacidad de batería. Así, la próxima vez que vayas a una tienda, podrás juzgar rápidamente si las “características” que te recomienda el comercial son necesidades reales o ventas forzadas. Por ejemplo:

\begin{itemize}
  \item Imprescindible: batería de más de 5000mAh, carga rápida de más de 80W, NFC, resistencia al agua IP68.
  \item Bonus: resumen con IA en el dispositivo, Wi-Fi 7, mensajes cortos de BeiDou.
  \item No necesario: pantalla plegable (estoy acostumbrado a barra recta), llamadas por satélite (no viajo lejos), carga rápida extrema de 240W.
\end{itemize}

\end{tcbpractice}

\begin{tcbsummary}{Resumen del capítulo}
\begin{itemize}
  \item Tres criterios para AI phone: modelo grande en el dispositivo, NPU TOPS, escenarios de AI Agent
  \item Estrategias de IA de los fabricantes: Apple privilegia privacidad e integración, Google modelos en el dispositivo de código abierto, Huawei/Xiaomi/OPPO/vivo cada uno con sus fortalezas
  \item Comunicación por satélite: mensajes cortos de BeiDou en gama media, llamadas Tiantong siguen siendo exclusivas de gama alta
  \item Pantallas plegables: hardware maduro, adaptación de apps aprox. 70\%, modo suspendido y pantalla dividida son experiencias únicas
  \item Estándar 6G en 2027, Wi-Fi 7 ya popularizado, LE Audio de Bluetooth empezando a aterrizar
  \item La batería de silicio-carbono hace que 5000-6000mAh sea estándar en gama alta, reduciendo drásticamente la ansiedad por la batería
\end{itemize}

\end{tcbsummary}
